% Thorium présent dans l'écorce terrestre
Le thorium \ch{_{90}^{237}Th} est radioactif et émetteur $\alpha$. La temps de
demi-vie $t_{1/2}$ du thorium 237 est égale à 18 jours. On dispose initialement
($t=0$) d'une source radioactive de thorium de masse $m_{0}=\qty{1}{\ug}$.
\begin{enumerate}
    \item Définir un noyau radioactif.~\textbf{(0,5pt)}
    \item Donner la composition du noyau \ch{_{90}^{237}Th}.~\textbf{(0,5pt)}
    \item Écrire l'équation de sa désintégration et préciser les lois de
          conservation utilisées sachant qu'elle produit du radium
          \ch{Ra}.~\textbf{(1pt)}
    \item Calculer la constante radioactive $\lambda$ du thorium
          237.~\textbf{(1pt)}
    \item Calculer la masse de thorium 237 restante à la date
          $t=\qty{36}{jours}$.~\textbf{(1pt)}
    \item À quelle date la masse de thorium 237 restante est-elle de
          $\qty{0.0156}{\ug}$~?~\textbf{(1pt)}
\end{enumerate}

